\section{Materials and methods}
\label{sec:methods}

\subsection{Data and analysis windows}
\label{sec:meth-data}

Weekly laboratory-confirmed influenza hospital admissions and the percentage of hospitals reporting them were taken from the NHSN Weekly Hospital Respiratory Data table for the 50 US states and the District of Columbia \citep{cdc2026nhsn}, using the version retrieved on 2026-10-20. The table contained 21641 jurisdiction-weeks from 2020-08-08 to 2026-10-10. Reporting was mandatory until 30 April 2024, voluntary from 1 May to 31 October 2024 and mandatory again from 1 November 2024, giving 6009, 420 and 4077 jurisdiction-weeks with a reporting scale of at least five in the three regimes.

The primary windows covered early epidemic growth in four seasons from 2022--23 to 2025--26. Each window contained the contiguous weeks from 1 October, or from 9 November 2024 when mandatory reporting resumed, in which national weekly growth was at least 1.05, and ended the week before the first national turning point. The windows ran from 2022-10-01 to 2022-11-26, 2023-10-07 to 2023-12-23, 2024-11-09 to 2024-12-28 and 2025-10-04 to 2025-12-27, giving 42 calendar weeks. This rule replaced a first version that ended each window at the season maximum and gave only four weeks in the double-peaked 2024--25 season; the revision was made on 2026-10-20 before any intercept was fitted on the new windows. For continuity, the three calendar windows of an earlier version of this analysis (2022-10-08 to 2022-11-26, 2023-10-14 to 2023-12-16 and 2024-11-09 to 2025-01-18) were analysed as a secondary specification. A peak window contained the national peak week of each season and the seven weeks after it. No zero-imputed counts entered any window.

\subsection{Weekly growth statistics}
\label{sec:meth-stats}

Two statistics of weekly growth were computed for each state and week. The weighted growth ratio is given by:
\begin{equation}
  G_t=\frac{C_{t+1}}{D_t},\qquad D_t=0.65\,C_t+0.25\,C_{t-1}+0.10\,C_{t-2},
  \label{eq:dt}
\end{equation}
where $C_t$ is the number of admissions in week $t$ and $D_t$ is the reporting scale, with $\sum_sw_s{}^2=0.495$. Only weeks with $D_t\ge5$ were used for $G$. The renewal-type ratio converts single-week growth into a reproduction number through the generation interval \citep{wallinga2007how}, and is given by:
\begin{equation}
  R_{A,t}=\left(1+r_t\theta\right)^{\alpha},\qquad r_t=\frac{1}{7}\log\frac{C_{t+1}+0.5}{C_t+0.5},
  \label{eq:ra}
\end{equation}
where $r_t$ is the daily growth rate, and $\alpha$ and $\theta$ are the shape and scale of a gamma generation interval with mean 3.6 days and standard deviation 1.6 days, the household serial interval estimated by \citet{cowling2009estimation}. Generation intervals with means of 2.2 and 4.3 days were used as sensitivity analyses. Because the scale of $R_A$ is the single-week count, it was computed only for weeks with $C_t\ge5$ as well as $D_t\ge5$. $R_A$ is a transformation of observed growth and not a posterior estimate of $R_t$ \citep{cori2013new}; both statistics are therefore apparent quantities whose relationship to $v_R$ was established by simulation (Section~\ref{sec:meth-inversion}).

\subsection{Expected loss and the balance scale}
\label{sec:meth-loss}

In routine surveillance the number of primary cases is unobserved, and secondary infections arise from the infection pressure $\Lambda_{t+1}=\sum_{s\ge1}w_sI_{t+1-s}$ \citep{fraser2007estimating, cori2013new}. Applying the law of total variance to $C_{t+1}$ given the history $\mathcal{F}_t$, with $m=\rho\E(\Lambda_{t+1}\mid\mathcal{F}_t)$, the relative prediction variance can be written as:
\begin{equation}
  \frac{\Var(C_{t+1}\mid\mathcal{F}_t)}{m^2}=v_R+\frac{A}{m}+\mathrm{CV}^2(\Lambda_{t+1}\mid\mathcal{F}_t)\,(\Rbar^2+v_R),
  \label{eq:info-decomp}
\end{equation}
where $\mathrm{CV}$ is the coefficient of variation. The last term is a further scale-invariant component that vanishes only when infection pressure is known. For a statistic $X_t$ with mean $\bar{X}$, the relative loss $L_t=[(X_t-\bar{X})/\bar{X}]^2$ was modelled as:
\begin{equation}
  \E(L_t\mid D_t)=a+\frac{b}{D_t},
  \label{eq:affine}
\end{equation}
where, if a week were a single generation and $D_t$ were free of noise, $a=v_R/\Rbar^2$, $b=A/\Rbar^2$ and $b/a=\mx$. Binomial noise in the denominator, with relative variance $(1-\rho)\sum_sw_s{}^2/m$, increases the slope to approximately:
\begin{equation}
  b\approx\frac{A+(1-\rho)(\Rbar^2+v_R)\sum_sw_s{}^2}{\Rbar^2}.
  \label{eq:b-delta}
\end{equation}
Theoretical ranges of $b/a$ for $G$ used $\Rbar$ equal to the mean of $G$ in each specification, $v_R\in\{0.04,0.08\}$, $k\in\{0.2,0.5\}$ \citep{lloydsmith2005superspreading} and an admission probability per infection $\rho\in\{0.0056,0.021\}$ \citep{tokars2018seasonal}, with and without Eq.~(\ref{eq:b-delta}). The uncorrected balance scale over these scenarios was 21.3 to 49.2 cases.

\subsection{Estimation and inference}
\label{sec:meth-inference}

Eq.~(\ref{eq:affine}) was fitted by feasible generalised least squares with weights $1/\hat{\mu}^2$, where $\hat{\mu}$ is the fitted mean from the previous iteration, iterated to convergence \citep{wooldridge2010econometric}. Observations from the same state are dependent over time, and observations from the same week share national influences. Uncertainty was therefore estimated by a two-way bootstrap that draws states and weeks independently with replacement and weights each observation by the product of the two draw counts \citep{cameron2008bootstrap, cameron2011robust}. Each specification used 2000 bootstrap replicates with its own seed (20261024 for $R_A$ and 20261124 for $G$ in the rule windows, 20261224 for $G$ and 20261324 for $R_A$ in the calendar windows), and uncertainty is reported as the point estimate followed by the 95\% percentile interval, for example 0.0150 (95\% CI: 0.0100, 0.0208).

The intercept was decomposed by refitting Eq.~(\ref{eq:affine}) twice. The first refit removed national-week effects by using deviations from the mean across states in each week, with a small-sample factor $n_w/(n_w-1)$. The second removed additive week and state-season effects estimated by alternating projections. The national-week component is the difference between the unadjusted and week-adjusted intercepts, and the state-season component is the further change on adding state-season effects. The affine model was compared with a power law $\E(L_t\mid D_t)=cD_t{}^{-\gamma}$ by a weighted contrast of residual sums of squares, calibrated by a residual bootstrap under the power law, and by leave-one-state-out and leave-one-season-out cross-validation. Power to detect a positive intercept was computed by simulating from the fitted affine model with intercepts up to the estimated value. Sensitivity analyses of $b/a$ trimmed the most influential 1\% of losses, omitted one season at a time, used Huber weights and removed national-week effects. Further analyses used the peak window, changed the generation interval and changed the denominator weights to a single week or three equal weeks.

\subsection{Reporting completeness}
\label{sec:meth-reporting}

Coverage-adjusted statistics were computed after dividing each count by the share of hospitals reporting in that state and week. The change in intercept, $\Delta a=a-a_{\mathrm{adj}}$, was tested for equivalence to zero by two one-sided tests at the 5\% level, using 90\% intervals from a paired two-way bootstrap \citep{schuirmann1987comparison}. The margins were set at the intercept that a relative reporting volatility $\tau_0$ would add under Proposition~\ref{prop:report-noise}, and are given by:
\begin{equation}
  \delta_G=\tau_0{}^2\left(1+\frac{v_R}{\Rbar^2}\right),\qquad
  \delta_{R_A}=\kappa_A{}^2\,\tau_0{}^2\left(1+\frac{v_R}{\Rbar^2}\right),
  \label{eq:tost}
\end{equation}
where $\kappa_A=\mathrm{d}\log R_A/\mathrm{d}\log g$ propagates a perturbation of weekly growth $g$ to $R_A$, $\Rbar$ is the sample mean of each statistic, $v_R=0.04$, and $\tau_0=0.10$ (primary) or 0.05. Because $\Rbar$ is estimated, the margins expressed as shares of the intercept carry bootstrap intervals. Power was computed from a normal approximation. The interquartile range of $\log G$ and the standard deviation of $G$ over May to October were compared between 2023, 2024 and 2025. As a second proxy, $R_A$ was computed from weekly positive influenza tests in clinical laboratories for 44 states with available data, and the share of the hospital intercept removed by regressing hospital $R_A$ on laboratory $R_A$ was estimated with 500 two-way bootstrap replicates.

\subsection{Seasons and pathogens}
\label{sec:meth-seasons}

Season-specific intercepts were pooled with a DerSimonian--Laird random-effects model \citep{dersimonian1986meta}, with a $t$ interval on $S-1$ degrees of freedom for the pooled intercept and a Q-profile interval for the between-season variance \citep{viechtbauer2007confidence}. Corollary~\ref{cor:season} gives the lower bound on the precision of any between-season variance \citep{searle1992variance}. Intercept ratios $q_1=a_{(1,0,0)}/a_{w}$ and $q_2=a_{(1/3,1/3,1/3)}/a_{w}$ compare a single-week and an equal three-week denominator with the weights of Eq.~(\ref{eq:dt}). They were computed for influenza in four windows and for COVID-19 and RSV admissions in two windows each. Null ranges came from simulated epidemic growth with reporting noise but no environmental variation, over three generation intervals, three admission delays and four noise levels. Each of the eight combinations was tested on five one-sided criteria ($q_1<1$, $q_2>1$, each ratio outside its null range and a positive excess over $v_R/\Rbar^2$), and the 40 tests were adjusted by the Holm and Benjamini--Hochberg procedures \citep{holm1979simple, benjamini1995controlling}.

\subsection{Simulation of apparent statistics}
\label{sec:meth-sim}

The closed forms of Section~\ref{sec:background} were checked by Monte Carlo with 20000 replications per point. In a first, generic design, apparent weekly statistics were simulated from a daily renewal model with $\Rbar=1.2$ and the generation interval above. Weekly $R_t$ was drawn from a gamma distribution with variance $v_R$, independently across weeks, offspring were negative binomial and admissions were binomial with probability $\rho_h$ after a gamma delay with mean 5.5 days and standard deviation 3.0 days, which was assumed. Each replicate contained 500 states with initial reporting scales log-uniform between 2 and 150, four warm-up weeks and six observation weeks. Seventeen scenarios crossed $v_R\in\{0.01,0.02,0.04,0.08,0.16\}$ with $k\in\{0.1,0.35,1\}$ at $\rho_h=0.01$ and added $\rho_h\in\{0.005,0.02\}$ at the baseline, each with 20 replicates. The noise-free limit of each intercept came from a mean-field run of 20000 states in which all count noise was replaced by expectations.

\subsection{Two-step inversion}
\label{sec:meth-inversion}

The empirical intercepts were mapped to $v_R$ in two steps. In the first step, a data-matched version of the simulator gave the expected fitted intercept of each statistic as a function of $v_R$. Each simulated panel had the cross-section and length of the analysis panel, with 51 states and the four rule windows of 9, 12, 8 and 13 observation weeks. The season mean $R_s$ was the renewal-type transform of the geometric-mean national weekly growth over each window, and initial admissions in each state equalled its observed admissions in the first window week. Weekly reproduction numbers were $R_{i,s,w}=R_s+\eta_{s,w}+\varepsilon_{i,s,w}$, where the national week shock $\eta$ has variance $0.45\,v_R$ and is shared by all states, the state shock $\varepsilon$ has variance $0.55\,v_R$, and values were truncated below at 0.05. The share of 0.45 was set between the national-week shares of 0.41 and 0.51 observed for the two statistics (Section~\ref{sec:res-multiscale}); matching the simulated intraclass correlation and national-week share to the data is therefore a consistency check, and setting the share to zero or 0.65 changed the implied variance for $R_A$ from 0.040 to 0.039 and 0.040 (Table~\ref{tab:s8}). Transmission was a daily renewal process with negative binomial offspring ($k=0.35$), admissions were binomial with probability $\rho_h=0.01$ after the delay above, and both statistics were computed and filtered exactly as in the data, with the analyst's generation interval of mean 3.6 days. The design crossed $v_R\in\{0.02,0.04,0.06,0.09\}$ with 200 panel replicates per value, and Eq.~(\ref{eq:affine}) was fitted to each replicate. A calibration line $\E(\hat{a})=\gamma_0+\gamma_1v_R$ was fitted by weighted least squares with weights equal to the inverse Monte Carlo variance of each mean, and the scale factor $\kappa=\E(\hat{a})/(v_R/\Rbar^2)$ was computed with $\Rbar$ the row-weighted true mean.

In the second step, the implied variance $v_R^{\mathrm{impl}}=(a-\gamma_0)/\gamma_1$ was computed for each statistic. Its interval combined 2000 joint two-way bootstrap draws of the two empirical intercepts, which share state and week multiplicities, with independent draws of $(\gamma_0,\gamma_1)$ from their Monte Carlo sampling distribution. Agreement between the two statistics was defined as non-rejection, that is, a 95\% interval of their difference that contains zero. For the primary design, equivalence was also tested by two one-sided tests on the ratio of the implied variances with bounds of 0.80 and 1.25, using the 90\% interval of the ratio draws. Sensitivity settings changed one element of the design at a time (Table~\ref{tab:s8}). They removed the national shock or raised its share to 0.65, added first-order autocorrelation of 0.5 across weeks, set the true generation interval to a mean of 2.2 or 4.3 days while the analyst kept 3.6 days, and set the delay mean to 3.5 or 8 days. One further setting lowered $k$ to 0.1, another replaced the transmission shock by a synchronous multiplicative shock to reported counts, and the generic design above served as a bridge. Four settings added at the final revision set the true generation interval to 2.9 days, the lower end of the interval of \citet{cowling2009estimation}, set the first-order autocorrelation to 0.25, and drew weekly admissions from a gamma--Poisson mixture with a variance-to-mean ratio of 1.5 or 2.0 (Table~\ref{tab:s10}). The independence of weekly shocks was checked with the pooled lag-1 autocorrelation of $x/\bar{x}-1$ within state-season blocks of each statistic, with an interval from a state-cluster bootstrap of 1000 draws, and compared with its distribution over simulated panels at $v_R=0.04$. The calibrated theoretical $b/a$ used the same design over $v_R\in\{0.04,0.08\}$, $k\in\{0.2,0.5\}$ and $\rho_h\in\{0.0056,0.021\}$, and the range of $b/a$ across single replicates gave the distribution expected from one panel of this size. Simulation seeds were derived from 20261020 by a seed sequence indexed by setting, value of $v_R$ and replicate block; the joint bootstrap and the calibration draws used 20261921 and 20261922.

